\documentclass{article}
\usepackage[ngerman]{babel}
\usepackage[utf8]{inputenc}
\usepackage[T1]{fontenc}

% --- Pakete für Diagramme und Code ---
\usepackage{pgfplots}
\pgfplotsset{compat=1.18}
\usepackage{tikz}
\usetikzlibrary{shapes.geometric, arrows, positioning}
\usepackage{listings}
\usepackage{xcolor}

% --- Farbdefinitionen für Listings ---
\definecolor{codegreen}{rgb}{0,0.6,0}
\definecolor{codegray}{rgb}{0.5,0.5,0.5}
\definecolor{codepurple}{rgb}{0.58,0,0.82}
\definecolor{backcolour}{rgb}{0.95,0.95,0.92}

% --- Stildefinitionen für Listings ---
\lstdefinestyle{latexstyle}{
    backgroundcolor=\color{backcolour},
    commentstyle=\color{codegreen},
    keywordstyle=\color{magenta},
    numberstyle=\tiny\color{codegray},
    stringstyle=\color{codepurple},
    basicstyle=\ttfamily\small,
    breakatwhitespace=false,
    breaklines=true,
    captionpos=b,
    keepspaces=true,
    numbers=left,
    numbersep=5pt,
    showspaces=false,
    showstringspaces=false,
    showtabs=false,
    tabsize=2,
    frame=single,
    rulecolor=\color{gray!30}
}
\lstdefinestyle{csvstyle}{
    backgroundcolor=\color{backcolour},
    basicstyle=\ttfamily\small,
    captionpos=b,
    numbers=none,
    frame=single,
    rulecolor=\color{gray!30},
    breaklines=true
}

% --- Beispiel-CSV-Datei ---
\begin{filecontents*}{daten.csv}
x,y
0,1
1,2
2,4
3,3
4,5
5,7
\end{filecontents*}

% --- TikZ Styles für Flowchart ---
\tikzstyle{startstop} = [rectangle, rounded corners, minimum width=3cm,
 minimum height=1cm, text centered, text width=3cm,
 draw=black, fill=red!20]
\tikzstyle{process} = [rectangle, minimum width=3cm, minimum height=1cm,
 text centered, text width=3cm, draw=black, fill=orange!20]
\tikzstyle{decision} = [diamond, minimum width=3cm, minimum height=1cm,
 text centered, draw=black, fill=green!20, aspect=1]
\tikzstyle{io} = [trapezium, trapezium left angle=70, trapezium right angle=110,
 minimum width=3.5cm, minimum height=1cm,
 text centered, text width=2cm, draw=black, fill=blue!20]
\tikzstyle{arrow} = [thick, ->, >=stealth]

\title{Beispieldokument: Liniendiagramm und Flowchart in \texttt{LaTeX}}
\author{Jakob – Abteilung Elektrotechnik, HTBLuVA Salzburg}
\date{\today}

\begin{document}
\maketitle

%%%%%%%%%%%%%%%%%%%%%%%%%%%%%%%%%%%%%%%%%%%%%%%%%%%%%%%%%%%%%%%%%%%%%%%%%%%%%%%
\section{Liniendiagramm aus CSV-Daten mit \texttt{pgfplots}}
%%%%%%%%%%%%%%%%%%%%%%%%%%%%%%%%%%%%%%%%%%%%%%%%%%%%%%%%%%%%%%%%%%%%%%%%%%%%%%%

\subsection{Einleitung}
In diesem Abschnitt wird gezeigt, wie mit \LaTeX{} und dem Paket \texttt{pgfplots}
ein Liniendiagramm aus einer externen CSV-Datei erstellt wird.  
Die Codeausschnitte werden farbig dargestellt und anschließend das fertige Diagramm gezeigt.

\subsection{Vorbereitung der Daten}
Die zu visualisierenden Daten liegen in einer CSV-Datei vor.  
Ihr Inhalt sieht wie folgt aus:

\begin{lstlisting}[
    style=csvstyle,
    caption={Inhalt der Datei \texttt{daten.csv}},
    label=lst:csv
]
x,y
0,1
1,2
2,4
3,3
4,5
5,7
\end{lstlisting}

\noindent
Die erste Zeile enthält die Spaltenüberschriften \texttt{x} und \texttt{y}.  
Jede weitere Zeile stellt ein Messpaar dar, das als Punkt im Diagramm visualisiert wird.

\subsection{Einlesen und Plotten der Daten}
Der zentrale Befehl zum Einlesen der CSV-Daten ist \texttt{\textbackslash addplot} in Kombination mit \texttt{table}.

\begin{lstlisting}[
    style=latexstyle,
    caption={Plotten der Daten aus einer CSV-Datei},
    label=lst:addplot
]
\addplot[
    blue,
    mark=*,
    line width=1pt
] table [
    col sep=comma,
    x=x,
    y=y
] {daten.csv};
\end{lstlisting}

\noindent
Hier wird die Datei \texttt{daten.csv} eingelesen, wobei die Spalte \texttt{x} auf die X-Achse
und die Spalte \texttt{y} auf die Y-Achse abgebildet wird.

\subsection{Vollständiges Diagramm}
Das folgende Listing zeigt die gesamte \texttt{axis}-Umgebung:

\begin{lstlisting}[
    style=latexstyle,
    caption={Komplette \texttt{axis}-Umgebung für das Liniendiagramm},
    label=lst:axis
]
\begin{tikzpicture}
  \begin{axis}[
    title={Liniendiagramm aus \texttt{daten.csv}},
    xlabel={X-Werte},
    ylabel={Y-Werte},
    grid=major,
    legend pos=north west,
    width=12cm,
    height=8cm,
    xmin=0, xmax=5,
    ymin=0, ymax=8,
    thick
  ]
  
  \addplot[
    blue,
    mark=*,
    line width=1pt
  ] table [
    col sep=comma,
    x=x,
    y=y
  ] {daten.csv};
  
  \legend{Messreihe}
  \end{axis}
\end{tikzpicture}
\end{lstlisting}

\noindent
Die Umgebung \texttt{axis} definiert das Koordinatensystem.  
Parameter wie \texttt{grid}, \texttt{width}, \texttt{height} oder \texttt{legend pos}
bestimmen Darstellung und Layout.

\subsection{Ergebnis}
Das resultierende Diagramm ist in Abbildung~\ref{fig:lineplot} dargestellt.

\begin{center}
\begin{tikzpicture}
  \begin{axis}[
    title={Liniendiagramm aus \texttt{daten.csv}},
    xlabel={X-Werte},
    ylabel={Y-Werte},
    grid=major,
    legend pos=north west,
    width=12cm,
    height=8cm,
    xmin=0, xmax=5,
    ymin=0, ymax=8,
    thick
  ]
  \addplot[
    blue,
    mark=*,
    line width=1pt
  ] table [
    col sep=comma,
    x=x,
    y=y
  ] {daten.csv};
  \legend{Messreihe}
  \end{axis}
\end{tikzpicture}
\end{center}
\label{fig:lineplot}

%%%%%%%%%%%%%%%%%%%%%%%%%%%%%%%%%%%%%%%%%%%%%%%%%%%%%%%%%%%%%%%%%%%%%%%%%%%%%%%
\section{Flowchart der Ball Balancing Platform (BBP) mit \texttt{TikZ}}
%%%%%%%%%%%%%%%%%%%%%%%%%%%%%%%%%%%%%%%%%%%%%%%%%%%%%%%%%%%%%%%%%%%%%%%%%%%%%%%

\subsection{Einleitung}
Im zweiten Teil wird gezeigt, wie mit dem Paket \texttt{TikZ} ein vollständiges
Flowchart für die Softwaresteuerung der Ball Balancing Platform (BBP) erstellt wird.
Das Diagramm bildet den Programmablauf des Regelkreises ab.

\subsection{Definition der TikZ-Styles}
Die grafischen Styles bestimmen Form, Farbe und Textausrichtung der Blöcke.

\begin{lstlisting}[
    style=latexstyle,
    caption={Definition der Styles für das Flowchart},
    label=lst:styles
]
\tikzstyle{startstop} = [rectangle, rounded corners, minimum width=3cm,
 minimum height=1cm, text centered, text width=3cm,
 draw=black, fill=red!20]
\tikzstyle{process} = [rectangle, minimum width=3cm, minimum height=1cm,
 text centered, text width=3cm, draw=black, fill=orange!20]
\tikzstyle{decision} = [diamond, minimum width=3cm, minimum height=1cm,
 text centered, draw=black, fill=green!20, aspect=1]
\tikzstyle{io} = [trapezium, trapezium left angle=70, trapezium right angle=110,
 minimum width=3.5cm, minimum height=1cm,
 text centered, text width=2cm, draw=black, fill=blue!20]
\tikzstyle{arrow} = [thick, ->, >=stealth]
\end{lstlisting}

\subsection{Erstellung der Knoten}
Jede Funktionseinheit wird durch einen \texttt{node}-Befehl dargestellt.  
Die Positionierung erfolgt relativ zu anderen Knoten.

\begin{lstlisting}[
    style=latexstyle,
    caption={Definition der Knoten für das Flowchart},
    label=lst:nodes
]
\node(start) [startstop] {Start};
\node(homing) [process, below of=start] {Homing};
\node(control_loop) [process, below of=homing] {Control Loop};
\node(camera) [io, below of=control_loop] {Read Camera};
\node(circle_detection) [process, below of=camera] {Circle Detection};
\node(read_mode) [io, below of=circle_detection] {Read Mode};
\node(mode1) [decision, below of=read_mode, yshift=-0.6cm] {Mode = 1?};
\node(mode2) [decision, below of=mode1, yshift=-1.2cm] {Mode = 2?};
\node(target_zero) [process, below of=mode2, yshift=-0.6cm] {Set Target = 0};
\node(read_joystick) [io, right of=mode1, xshift=6cm] {Read Joystick};
\node(target_circle) [process, right of=mode2, xshift=2cm] {Set Target = Circle()};
\node(calc_target) [process, right of=mode2, xshift=6cm] {Compute Target Value};
\node(pid) [process, below of=target_zero] {PID Control};
\node(invkin) [process, below of=pid] {Inverse Kinematics};
\node(motor_control) [io, below of=invkin] {Drive Motors};
\node(exit) [decision, below of=motor_control, yshift=-0.5cm] {Exit?};
\node(cleanup) [process, below of=exit, yshift=-0.5cm] {Cleanup};
\node(stop) [startstop, below of=cleanup] {Stop};
\end{lstlisting}

\subsection{Verbindungen und Pfeile}
Die logischen Verbindungen werden mit \texttt{\textbackslash draw[arrow]} erstellt.

\begin{lstlisting}[
    style=latexstyle,
    caption={Verbindungen der Blöcke mit Beschriftungen},
    label=lst:arrows
]
\draw[arrow] (mode1) -- (mode2)         node[anchor=west, midway] {No};
\draw[arrow] (mode1) -- (read_joystick) node[anchor=south, midway] {Yes};
\draw[arrow] (mode2) -- (target_zero)   node[anchor=west, midway] {No};
\draw[arrow] (mode2) -- (target_circle) node[anchor=south, midway] {Yes};
\draw[arrow] (target_circle) |- (pid);  % Knickpfad
\draw[arrow] (exit.east) -- node[anchor=south] {No} ++(10,0) |- (control_loop.east);
\end{lstlisting}

\subsection{Ergebnis}
Das fertige Flowchart ist in Abbildung~\ref{fig:flowchart} dargestellt.

\begin{center}
\begin{tikzpicture}[node distance=2cm, scale=.55, transform shape]
\node(start) [startstop] {Start};
\node(homing) [process, below of=start] {Homing};
\node(control_loop) [process, below of=homing] {Control Loop};
\node(camera) [io, below of=control_loop] {Read Camera};
\node(circle_detection) [process, below of=camera] {Circle Detection};
\node(read_mode) [io, below of=circle_detection] {Read Mode};
\node(mode1) [decision, below of=read_mode, yshift=-0.6cm] {Mode = 1?};
\node(mode2) [decision, below of=mode1, yshift=-1.2cm] {Mode = 2?};
\node(target_zero) [process, below of=mode2, yshift=-0.6cm] {Set Target = 0};
\node(read_joystick) [io, right of=mode1, xshift=6cm] {Read Joystick};
\node(target_circle) [process, right of=mode2, xshift=2cm] {Set Target = Circle()};
\node(calc_target) [process, right of=mode2, xshift=6cm] {Compute Target Value};
\node(pid) [process, below of=target_zero] {PID Control};
\node(invkin) [process, below of=pid] {Inverse Kinematics};
\node(motor_control) [io, below of=invkin] {Drive Motors};
\node(exit) [decision, below of=motor_control, yshift=-0.5cm] {Exit?};
\node(cleanup) [process, below of=exit, yshift=-0.5cm] {Cleanup};
\node(stop) [startstop, below of=cleanup] {Stop};
\draw[arrow] (start) -- (homing);
\draw[arrow] (homing) -- (control_loop);
\draw[arrow] (control_loop) -- (camera);
\draw[arrow] (camera) -- (circle_detection);
\draw[arrow] (circle_detection) -- (read_mode);
\draw[arrow] (read_mode) -- (mode1);
\draw[arrow] (mode1) -- (mode2) node[anchor=west, midway] {No};
\draw[arrow] (mode1) -- (read_joystick) node[anchor=south, midway] {Yes};
\draw[arrow] (mode2) -- (target_zero) node[anchor=west, midway] {No};
\draw[arrow] (mode2) -- (target_circle) node[anchor=south, midway] {Yes};
\draw[arrow] (read_joystick) -- (calc_target);
\draw[arrow] (target_circle) |- (pid);
\draw[arrow] (calc_target) |- (pid);
\draw[arrow] (target_zero) -- (pid);
\draw[arrow] (pid) -- (invkin);
\draw[arrow] (invkin) -- (motor_control);
\draw[arrow] (motor_control) -- (exit);
\draw[arrow] (exit) -- node[anchor=west] {Yes} (cleanup);
\draw[arrow] (cleanup) -- (stop);
\draw[arrow] (exit.east) -- node[anchor=south] {No} ++(10,0) |- (control_loop.east);
\end{tikzpicture}
\end{center}
\label{fig:flowchart}

\end{document}
